% GENERATED FILE. Edit canonical lessons, then run npm run generate:books.
% canonical-source-hash: fnv1a64:6fb373d82790d76a
% canonical-lessons: IT-C529-doppiare, IT-C529-doppiaggio, IT-C529-sottotitoli, IT-C529-colonna-sonora, IT-C529-trailer

\chapter{In lingua originale}
\label{ch:it-in-lingua-originale}
\hlchaptermodality{529}

\begin{chapteropening}
\textbf{By the end of this chapter:} \emph{I can talk about watching films dubbed or with subtitles.}

Complete the last lesson of chapter 529: I can talk about watching films dubbed or with subtitles.
\end{chapteropening}

\section[doppiare]{doppiare — to dub}

\label{lesson:IT-C529-doppiare}



\begin{quote}
Before the new one: say the Italian for spoiler, then the Italian for remote control.
\end{quote}

\subsection*{You'll want to know: doppiare}
\textbf{doppiare} — \textquotedblleft{}to dub\textquotedblright{}.

\begin{quote}
Quasi tutti i film stranieri sono doppiati.
\end{quote}

\begin{quote}
\emph{Almost all foreign films are dubbed.}
\end{quote}

Also to lap another runner, or to sail round a cape.

\begin{cousinweb}
\emph{doppiare} is inherited from Late Latin \emph{dūplāre} \textquotedblleft{}to double\textquotedblright{} (EtymDB), from \emph{duplus} \textquotedblleft{}double\textquotedblright{}.
\end{cousinweb}

\begin{grammarlens}[title={What you've built}]
One more word for telling, explaining and describing at B1.
\end{grammarlens}

\subsection*{Your turn}
\begin{itemize}
\raggedright
  \item \emph{Say it:} \emph{doppiare}
  \item \emph{Say it:} \emph{doppiare}, once more
  \item \emph{Say it:} a sentence of your own with \emph{doppiare}
  \item \emph{Recall:} say the Italian for spoiler, then the Italian for remote control, then say \emph{doppiare} again
\end{itemize}

\begin{tcolorbox}[breakable,colback=teal!4,colframe=teal!35!black,title={Before you move on}]
What does doppiare mean? (\textquotedblleft{}To dub\textquotedblright{}.) Say it once more.
\end{tcolorbox}
\section[il doppiaggio]{il doppiaggio — dubbing}

\label{lesson:IT-C529-doppiaggio}



\begin{quote}
Before the new one: say the Italian for remote control, then the Italian for to dub.
\end{quote}

\subsection*{You'll want to know: il doppiaggio}
\textbf{il doppiaggio} — \textquotedblleft{}dubbing\textquotedblright{}.

\begin{quote}
Il doppiaggio di questo film è ottimo.
\end{quote}

\begin{quote}
\emph{The dubbing of this film is excellent.}
\end{quote}

The person is \emph{il doppiatore}.

\begin{cousinweb}
\emph{doppiaggio} is built on \emph{doppiare} \textquotedblleft{}to dub\textquotedblright{}, from Late Latin \emph{dūplāre} \textquotedblleft{}to double\textquotedblright{} (EtymDB), with \emph{-aggio} for the work.
\end{cousinweb}

\begin{grammarlens}[title={What you've built}]
One more word for telling, explaining and describing at B1.
\end{grammarlens}

\subsection*{Your turn}
\begin{itemize}
\raggedright
  \item \emph{Say it:} \emph{il doppiaggio}
  \item \emph{Say it:} \emph{il doppiaggio}, once more
  \item \emph{Say it:} a sentence of your own with \emph{il doppiaggio}
  \item \emph{Recall:} say the Italian for remote control, then the Italian for to dub, then say \emph{il doppiaggio} again
\end{itemize}

\begin{tcolorbox}[breakable,colback=teal!4,colframe=teal!35!black,title={Before you move on}]
What does il doppiaggio mean? (\textquotedblleft{}Dubbing\textquotedblright{}.) Say it once more.
\end{tcolorbox}
\section[i sottotitoli]{i sottotitoli — subtitles}

\label{lesson:IT-C529-sottotitoli}



\begin{quote}
Before the new one: say the Italian for to dub, then the Italian for dubbing.
\end{quote}

\subsection*{You'll want to know: i sottotitoli}
\textbf{i sottotitoli} — \textquotedblleft{}subtitles\textquotedblright{}.

\begin{quote}
Guardo i film in inglese con i sottotitoli.
\end{quote}

\begin{quote}
\emph{I watch films in English with subtitles.}
\end{quote}

The singular \emph{il sottotitolo} is also a subheading.

\begin{cousinweb}
\emph{sottotitoli} is built from \emph{sotto} \textquotedblleft{}under\textquotedblright{}, from Latin \emph{subtus}, and \emph{titolo} \textquotedblleft{}title\textquotedblright{}, from Latin \emph{titulus} (Lewis \& Short).
\end{cousinweb}

\begin{grammarlens}[title={What you've built}]
One more word for telling, explaining and describing at B1.
\end{grammarlens}

\subsection*{Your turn}
\begin{itemize}
\raggedright
  \item \emph{Say it:} \emph{i sottotitoli}
  \item \emph{Say it:} \emph{i sottotitoli}, once more
  \item \emph{Say it:} a sentence of your own with \emph{i sottotitoli}
  \item \emph{Recall:} say the Italian for to dub, then the Italian for dubbing, then say \emph{i sottotitoli} again
\end{itemize}

\begin{tcolorbox}[breakable,colback=teal!4,colframe=teal!35!black,title={Before you move on}]
What does i sottotitoli mean? (\textquotedblleft{}Subtitles\textquotedblright{}.) Say it once more.
\end{tcolorbox}
\section[la colonna sonora]{la colonna sonora — soundtrack}

\label{lesson:IT-C529-colonna-sonora}



\begin{quote}
Before the new one: say the Italian for dubbing, then the Italian for subtitles.
\end{quote}

\subsection*{You'll want to know: la colonna sonora}
\textbf{la colonna sonora} — \textquotedblleft{}soundtrack\textquotedblright{}.

\begin{quote}
La colonna sonora del film è bellissima.
\end{quote}

\begin{quote}
\emph{The film's soundtrack is beautiful.}
\end{quote}

Literally \textquotedblleft{}the sound strip\textquotedblright{}: on old film, the strip that carried the sound.

\begin{cousinweb}
\emph{colonna sonora} is built from \emph{colonna}, inherited from Latin \emph{columna} \textquotedblleft{}a pillar, a post\textquotedblright{} (EtymDB), and \emph{sonoro} \textquotedblleft{}of sound\textquotedblright{}, from Latin \emph{sonōrus} (Lewis \& Short).
\end{cousinweb}

\begin{grammarlens}[title={What you've built}]
One more word for telling, explaining and describing at B1.
\end{grammarlens}

\subsection*{Your turn}
\begin{itemize}
\raggedright
  \item \emph{Say it:} \emph{la colonna sonora}
  \item \emph{Say it:} \emph{la colonna sonora}, once more
  \item \emph{Say it:} a sentence of your own with \emph{la colonna sonora}
  \item \emph{Recall:} say the Italian for dubbing, then the Italian for subtitles, then say \emph{la colonna sonora} again
\end{itemize}

\begin{tcolorbox}[breakable,colback=teal!4,colframe=teal!35!black,title={Before you move on}]
What does la colonna sonora mean? (\textquotedblleft{}Soundtrack\textquotedblright{}.) Say it once more.
\end{tcolorbox}
\section[il trailer]{il trailer — trailer (of a film)}

\label{lesson:IT-C529-trailer}



\begin{quote}
Before the new one: say the Italian for subtitles, then the Italian for soundtrack.
\end{quote}

\subsection*{You'll want to know: il trailer}
\textbf{il trailer} — \textquotedblleft{}trailer (of a film)\textquotedblright{}.

\begin{quote}
Il trailer mi è piaciuto, andiamo a vederlo.
\end{quote}

\begin{quote}
\emph{I liked the trailer, let us go and see it.}
\end{quote}

From English; said as in English.

\begin{cousinweb}
\emph{trailer} was borrowed from English \emph{trailer}, built on \emph{trail} \textquotedblleft{}to drag behind\textquotedblright{}, usually traced to Latin \emph{trāgula} \textquotedblleft{}a drag-net\textquotedblright{} (Lewis \& Short).
\end{cousinweb}

\begin{grammarlens}[title={What you've built}]
One more word for telling, explaining and describing at B1.
\end{grammarlens}

\subsection*{Your turn}
\begin{itemize}
\raggedright
  \item \emph{Say it:} \emph{il trailer}
  \item \emph{Say it:} \emph{il trailer}, once more
  \item \emph{Say it:} a sentence of your own with \emph{il trailer}
  \item \emph{Recall:} say the Italian for subtitles, then the Italian for soundtrack, then say \emph{il trailer} again
\end{itemize}

\begin{tcolorbox}[breakable,colback=teal!4,colframe=teal!35!black,title={Before you move on}]
What does il trailer mean? (\textquotedblleft{}Trailer (of a film)\textquotedblright{}.) Say it once more.
\end{tcolorbox}
